% Lösungen Einheit 1 – Terme aufstellen und berechnen (zusammenbau v1.0)
\kbabschnitt{Terme aufstellen und berechnen}

\erg{27}{Vierfaches $4x$, Summe $4x + 9$, halbiert mit Klammer: $(4x + 9) : 2$ (richtig $\rightarrow$ Nr.~29–30 überspringen)}
\erg{28}{a) $4 \cdot 5 - 3 = 17$ \quad b) $2 \cdot (-3) + 9 = 3$ \quad c) $10 - 3 \cdot (-2) = 16$}
\erg{29}{a) $4x - 3$ \quad b) $4x - 5$ \quad c) $4x - 6$ \quad d) $3x + 8$ \quad e) $5x$ (auch $5 \cdot x$) \quad f) $x + 9 + x + 9 = 2x + 18$ \quad g) $2 \cdot (3x + 2)$}
\erg{30}{a) Dreifaches $3x$, Differenz $3x - 7$, verdoppelt mit Klammer: $2 \cdot (3x - 7)$ \quad b) Sechsfaches $6x$, dann um 2 vermindert: $6x - 2$}
\erg{31}{a) Zum Beispiel: Ein Kinobesuch kostet 15 € Eintritt, dazu kommen $x$ Getränke zu je 3 €; $x$ ist die Zahl der Getränke. Einsetzen: $15 + 3 \cdot 2 = 21$, der Besuch kostet 21 €. \quad b) Zum Beispiel: Lara hat 20 € und kauft $x$ Hefte zu je 2 €; $x$ ist die Zahl der Hefte. Für $x = 4$: $20 - 2 \cdot 4 = 12$, ihr bleiben 12 €. \quad c) Zum Beispiel: Ein Fahrradverleih nimmt 10 € Grundgebühr und 4 € je Stunde; $x$ ist die Zahl der Stunden. Für $x = 3$: $10 + 4 \cdot 3 = 22$, drei Stunden kosten 22 €.}
\erg{32}{a) Richtig: $x - 9$. Mia hat die Reihenfolge vertauscht; „vermindert um 9“ heißt, die 9 wird von $x$ abgezogen. \quad b) Die Klammer sorgt dafür, dass die ganze Summe verdoppelt wird, nicht nur $x$; Probe mit $x = 3$: $2 \cdot (3 + 5) = 16$, aber $2 \cdot 3 + 5 = 11$ – die Terme haben verschiedene Werte. \quad c) $4 \cdot (x + 6)$ \quad d) Term aufstellen: $11 + 0{,}30n$; einsetzen: $11 + 0{,}30 \cdot 150 = 56$ €. Ein Monat mit 150 Kilowattstunden kostet 56 €.}
